\subsection{NN-G025: the exact 4520-pair optimum for proportional cells}
\paragraph{Statement.}
The exact minimum is $4520$ cancelling pairs in the closed
rank-three table class with arbitrary proportional effective
cell polynomials $\delta_1,\delta_2\in\mathbb Z[D]$,
arbitrary-degree $W\in\mathbb Z[D]$, positive rational
constant ratios $G_i/(W\delta_i)$, and arbitrary rational
auxiliary functions $\Lambda_i$. All actual coefficients
must obey the usual integrality, domination, effectivity
and parity conditions. Proportional means
$\delta_1/\delta_2\in\mathbb Q_{>0}$; the contents need
not agree and the common primitive polynomial has no
degree restriction.

An attainer has $3532$ bottom-cell cancelling pairs.
This enlarges NN-G024's equal-cell class, whose exact
minimum $5932$ remains valid. The new global upper
bound is $4520$; global minimality without the
proportional-cell and constant-ratio restrictions is
not asserted.

\paragraph{The attainer and its exact score.}
Let $f=(D-1)(D-2)$. Set
\[
 \begin{aligned}
 \delta_1&=f,&\delta_2&=3f,&W&=3D,\\
 G_1&=6Df,&G_2&=5Df,&S&=10Df^2,\\
 U_1&=3D,&U_2&=D,\\
 \Lambda_1&=6D^2-25D+20,&
 \Lambda_2&=(D^2-2D-4)/3,\\
 V_1&=15D-12,&V_2&=2D^2+D+2.
 \end{aligned}                                                     \tag{G025.1}
\]
Here $G_i=k_iW\delta_i$ with $(k_1,k_2)=(2,5/9)$.
Direct substitution gives
\[
 \frac{G_1G_2}{S}=W,\qquad
 \frac{G_i(G_i-D\delta_i)}{S}=U_i,\qquad
 \frac{G_i(DG_i-\delta_i\Lambda_i)}{S}=V_i.
\]
The injective two-matrix-block/two-scalar-block construction
of NN-G011 works with the two different $\delta_i$ and
these nonzero $\Lambda_i(D^2-\Lambda_i)$, proving
associativity.

The changed effective base coefficients, in symmetric
nonnegative-index Laurent lists, are
\[
\begin{array}{c|l|l}
 &\text{dimension coefficient}&\text{absolute trace target}\\\hline
 \delta_1&(2,1,1)&(0,1,1)\\
 \delta_2&(6,3,3)&(0,1,1)\\
 \delta_1\Lambda_1&(22,9,17,5,6)&(2,1,1,1)\\
 \delta_2\Lambda_2&(2,5,2,3,1)&(2,3,2,1,1)\\
 G_1&(36,30,18,6)&(16,12,4)\\
 G_2&(30,25,15,5)&(14,11,5,1)\\
 S&(280,250,180,100,40,10)&(144,120,68,24,4)\\
 V_1&(18,15)&(6,5)\\
 V_2&(16,9,2)&(6,3,2)
\end{array}                                                        \tag{G025.2}
\]
Also $W=U_1=3D$ and $U_2=D$. These lists prove effectivity,
domination and parity for all base coefficients; powers
of $D$ propagate them to the full table exactly as in
NN-G024. No actual coefficient has a denominator.
The independent integer-Laurent checker verifies all
$132$ positions and all $1000$ associativity triples.

At $D=4$ the bottom mass is
\[
 1440+20(144+120)+200\cdot12+20(12+4+48+38)=11160.
\]
The cell mass is
$26(6+18)+14(6\cdot16+6\cdot4)=2304$.
Thus $T=13464$ and
\[
 (13464-4424)/2=4520,\qquad(11160-4096)/2=3532.
 \tag{G025.3}
\]

\paragraph{Normalization and complete necessary bounds.}
Write $\delta_i=c_i f$ with $f$ primitive effective and
$c_i$ positive odd integers, and $W=gW_0$ with $g$
positive odd and $W_0$ primitive effective. Put
$d=f(4)$ and $w=W_0(4)$. The same Gauss and parity
arguments as NN-G024 give positive integers $a_1$ even
and $a_2$ odd, with $g\mid a_1a_2$, and
\[
 \begin{aligned}
 G_i&=a_iW_0f,&S&=(a_1a_2/g)W_0f^2,\\
 U_i&=g(a_iW_0-c_iD)/a_j,&
 V_i&=g(a_iDW_0-\Psi_i)/a_j,
 \end{aligned}                                                      \tag{G025.4}
\]
where $\Psi_i=c_i\Lambda_i\in\mathbb Z[D]$.
Integrality of $\Psi_i$ follows from polynomiality of
$V_i$ and the primitive product $f\Psi_i$, exactly
as before. It does not require $c_1=c_2$.
The actual cell-product masses are $dp_i$ with
$p_i=\Psi_i(4)$, and $dp_1\ge8$, $dp_2\ge16$.
The exact total mass is
\[
 \begin{split}
 T={}&(a_1a_2/g)wd^2+20wd(a_1+a_2)+200gw\\
 &+20\sum_i\{U_i(4)+V_i(4)\}
       +26(c_1+c_2)d+14d(p_1+p_2).
 \end{split}                                                        \tag{G025.5}
\]
For a proposed improvement on $4520$, the complete
bounds (G024.6) still hold, since $c_1+c_2\ge2$.
Moreover
\[
 T\ge3024+104(c_1+c_2).
\]
Using even the weaker cap $T<16288$ gives
$c_1+c_2\le126$, so each positive odd $c_i\le125$.
This generous bound suffices; no degree bound on $f$
or on outer Laurent support locations is being used.

Use $b_0,b_1$ for the zero and first coefficients of
$W_0$, and $M,h$ as in (G024.7), so
$w=b_0+2b_1+2Mh$.
The complete scalar ranges are those of NN-G024,
with the one content variable replaced by
$1\le c_i\le125$, both odd. For each $i$ impose
\[
 g(a_i b-c_i t)/a_j\in t+2\mathbb N
 \quad\text{for }(b,t)=(b_0,2),(b_1,1).               \tag{G025.6}
\]
The bound from the first coefficient also permits
$c_i\le a_i b_1$ without removing any actual case.
The outer-coefficient divisibility, primitive-content
necessary tests, $S(4)\ge576$, and $G_i(4)\ge48$ are
unchanged. The intervals and endpoint rule (G024.8)
for $p_i$ are unchanged as well. When $h=0$, the
leading-coefficient test (G024.9) is still necessary.

For fixed other parameters, the coefficient of $c_i$
in (G025.5) is $26d-80g/a_j$. Thus minimize the
content contribution at the least admissible $c_i$
when that coefficient is nonnegative, and at the
greatest otherwise. The two contents can be minimized
independently for a necessary relaxation: ignoring
their joint upper bound only enlarges the possible
set. Their individual bounds are imposed. An endpoint
row is a lower bound for its whole class of content
choices, not a claim that all those choices coincide.

The exact integer-only script
\texttt{verify\_proportional\_cells\_optimum.py}
uses precisely these ranges, congruences and endpoint
rules. It leaves just four rows with relaxed score
below $4520$:
\[
\begin{array}{c|c|c|c|c|r}
(a_1,a_2,g)&(c_1,c_2)_{\rm endpoint}&d&(b_0,b_1,h)&
 (p_1,p_2)_{\rm endpoint}&\text{relaxed score}\\\hline
(6,5,3)&(1,3)&4&(2,1,0)&(6,4)&2636\\
(4,9,3)&(1,5)&4&(2,1,0)&(4,8)&3480\\
(6,9,3)&(3,3)&4&(2,1,0)&(6,4)&3716\\
(6,5,3)&(1,3)&6&(2,1,0)&(6,4)&4160
\end{array}                                                        \tag{G025.7}
\]
All have $W_0=D$. In each of the first three,
$f$ has mass four and is consequently $D(D-3)$.
Its product with $D$ has central coefficient two.
Thus their $G_1$ central coefficients are respectively
$12,8,12$, all below the trace target $16$.
These exclusions are independent of the choice of
contents within the row.

\paragraph{The last row and an exact residue obstruction.}
In the remaining row the effective parity-correct
primitive polynomial of mass six is uniquely
$f=(D-1)(D-2)$: relative to the mass-four polynomial,
the only possible addition of mass two is a central
coefficient two. Condition (G025.6) at the first
coefficient forces $c_1=1$ and $c_2=3$. Indeed
$3(6-c_1)/5$ must be a positive odd integer, giving
$c_1=1$, and $(5-c_2)/2$ must be a positive odd
integer, giving $c_2=3$.

Since $W_0=D$, effectivity of $f\Psi_i$ and $V_i$
forces $\deg\Psi_i\le2$ by their opposite leading
signs above degree two. The congruence for $V_1$ is
\[
 \Psi_1\equiv6D^2-5(5D-4)\pmod{10}.
\]
Writing $\Psi_1=L_0+L_1(q+q^{-1})+L_2(q^2+q^{-2})$
gives
\[
 L_0\equiv6,\quad L_1\equiv9,\quad L_2\equiv6
       \pmod{10}.                                      \tag{G025.8}
\]
The highest product coefficient forces $L_2\ge0$,
and $V_1=(3/5)(6D^2-\Psi_1)$ forces $L_2\le6$.
Hence $L_2=6$. The mass $p_1=L_0+2L_1+12$ is
congruent to six modulo ten; $6p_1\ge8$ makes
six its smallest possible positive value.

Suppose $p_1=6$. Then $L_0=-6-2L_1$.
The central coefficient of $f\Psi_1$ is
$2L_0+2L_1+12=-2L_1$ and must be at least two,
so $L_1\le-1$. Its coefficient at $q^3$ is
$L_1+6$ and must be at least one, so $L_1\ge-5$.
Together with (G025.8) this forces $L_1=-1$ and
$L_0=-4$. But its coefficient at $q$ would then be
$L_0+3L_1+6=-1$, below the target one.
Thus $p_1\ne6$, and consequently $p_1\ge16$.
The second mass is a multiple of four and satisfies
$6p_2\ge16$, so $p_2\ge4$.

Substitution in the exact mass identity for this row gives
\[
 T=12016+72p_1+74p_2
   \ge12016+72\cdot16+74\cdot4=13464.
\]
It cannot improve on $4520$. All four rows of the
complete necessary relaxation have now been excluded
below that score, and (G025.1) attains it. This proves
the exact minimum in the stated class.\hfill$\square$

\paragraph{Geometric meaning and residual scope.}
The two cells have the same primitive grading character
and different positive multiplicities. Their actual
coefficient objects still lie in the effective graded
sign-representation semiring, and NN-G010 supplies the
same finite-kernel monoidal realization. The precise
classical dictionaries are SGA 3, I, 4.7.2--4.7.3 and
archive cote 2, scan 17. Rational $\Lambda_2$ remains
an auxiliary factor, not a fractional effective object.
Different contents of the two cell characters do not
alter the supplied trace algebra whose effective lift
is being sought; they change the chosen dimension data.

Among the NN-G010 kernel realizations whose two cell characters
are positive integral multiples of the same primitive character
and whose specified ratios are constant, the least total
half-excess of dimension over absolute signed trace is $4520$,
with $3532$ in the bottom-cell coefficients for (G025.1).
These numerical invariants refer to the chosen table and
generators; they are not invariants of every categorification
of the abstract trace algebra.

The exact common-cell minimum and the exact proportional-cell
minimum differ, so equality of cell characters is a
substantive optimization restriction, not a consequence
of the trace constraints. A smaller score must now use
nonproportional cell polynomials, nonconstant ratios,
or another table shape. No canonicality, generic
rank-four positivity, PH0/PH1 or nonstandard purity
conclusion is inferred from this noncanonical optimum.

